\section{已核事实与译本勘误}
\label{app:errata}

本附录只收“已经核过、可以直接引用”的事实；凡未逐条对照原刊页图者，一律标
\texttt{[待核]}，不混入下文。每条给页码或文件行号，便于研讨班现场复查。
附录 B 已列过其中三条的简版；本附录给完整清单、证据与本次新发现，两者若有出入，
以本附录为准（同一事实只在此处给全，正文与附录 B 只引用、不再重复论证）。

\subsection{可以直接引用的已核事实}

\begin{enumerate}[leftmargin=1.8em,itemsep=2pt]
\item \textbf{页数}：原刊 49 页，印刷页 351--399；中译本是 40 页重排，二者不能按页套。
      页边注“原文 p.$N$”只对原文页图成立（\src{351--399}）。
\item \textbf{引言没有“解决第二问”}：原刊引言只说“这两个问题可以粗略地描述如下”
      \src{351}；译本卷首导语框里“本文解决第二个问题”一类话是\textbf{译者所加}的导读，
      不是 Witten 原文。引用时务必区分。
\item \textbf{原始出处}：本文是 1988 年 7 月 Swansea IAMP 大会报告的扩充版
      （原刊 p.~351 脚注 *：“An expanded version of a lecture at the IAMP Congress,
      Swansea, July, 1988”）；收稿 1988-09-12，修改稿 1988-09-27，A.~Jaffe 送审。
\item \textbf{(1.2) 没有 $1/4$、没有负号}：原刊 p.~353 照印下式（无 $1/4$、无负号）：
      \begin{equation*}
      \mathcal L_0=\int_M\sqrt g\,g^{ik}g^{jl}\Tr(F_{ij}F_{kl}),
      \end{equation*}
      早期材料说“漏
      $1/4$”，2026-11-13 由 S1 证据官对照 \texttt{pages/page-03.png} 判\textbf{否}：
      与 $-\frac14\int\sqrt g\,\Tr F^2$ 的差别来自脚注 1 的 $\Tr$ 归一化与反厄米约定，
      不是笔误。\textbf{不要按 $-\frac14\int\Tr F^2$ 去“补”系数。}
\item \textbf{(3.2)/(3.3) 误引三处}：原刊 p.~367 末与 p.~368 共三处把高斯定律约束
      “the constraints (3.2)”误引，应为 \textbf{(3.3)}（$\varepsilon^{ij}F^a_{ij}=0$）；
      (3.2) 是泊松括号。译本照录原文并在页边注说明（\texttt{s04-canonical.tex:125,137,142}
      对应的页图：\texttt{page-17.txt:44}、\texttt{page-18.txt:6}、\texttt{page-18.txt:15}）。
      此错误只在 §3.2，不波及 §1（§1 全文不引 §3 的公式）。
\item \textbf{“Proof of Verlinde's Conjecture”的位置}：它是 p.~391 段首的斜体小标题
      （run-in heading），\textbf{不是} §4.5 的节标题；§4.5 的标题是
      “Some Concrete Surgeries”（p.~388）。引用时不要写成“§4.5 证明了 Verlinde 猜想”。
\item \textbf{Verlinde 猜想的证明状态}：Witten 在 p.~391 明说该猜想“已由 Moore 与
      Seiberg [13] 证明”，他给出流代数情形的“新证明”（物理推导）；
      数学上的完整证明要等 Faltings 1994（Schottenloher 中译 p.~177 明说
      “物理结论在数学中将被看作猜想”）。
\item \textbf{参考文献 39 条}，p.~397--399；引用编号 [n] 时以原刊编号为准。
\item \textbf{四本教材的页码约定}：河野中译 PDF $=$ 书页 $+14$；Moore 中译的页脚页
      $=$ PDF 页（不要信 chunks 头注：\texttt{c12.tex} 头注与页脚差约 5 页，
      \texttt{c29.tex} 头注超过全书 234 页）；上野 PDF $=$ 书页 $+12$；
      Schottenloher PDF $=$ 书页 $+12$。
\item \textbf{上野的 $S$ 矩阵在 p.~111}（\texttt{src/b12.tex:570} 起），(5.29)--(5.30)
      同页区；p.~103 属 §5.3，不是 $S$ 矩阵。
\item \textbf{河野的“定理 2.16”不是 Verlinde}：p.~60--61 的定理 2.16 是用 Dehn 手术
      定义三维流形 Witten 不变量（Kirby 不变性）；Verlinde 公式在命题 1.22（p.~37）。
      引用时别把两者当“同一公式的两种表述”。
\item \textbf{Schottenloher 的 Verlinde 在第 11 章}（书 p.~172--189，PDF 184--201），
      不在第 6 章；第 6 章是 Virasoro 代数表示论（Verma 模、Kac 行列式）。
\end{enumerate}

\subsection{读前必看的三个坑}

\begin{enumerate}[leftmargin=1.8em,itemsep=2pt]
\item \textbf{(3.2) 应为 (3.3)}：见上第 5 条。读 §3.2 时把三处“constraints (3.2)”
      在心里改成 (3.3)，否则会以为约束是泊松括号。
\item \textbf{不要给 (1.2) 补 $1/4$}：见上第 4 条。所有“Witten 的 $-\frac14\int\Tr F^2$”
      都是把脚注 1 的 $\Tr$ 归一换掉之后的重写，不是原式。
\item \textbf{“本文没有半形式修正”是译本批注}：Witten 原文没有出现 half-form 一词；
      “严格量子化还需 $K^{1/2}$、一般认为正是 $k\to k+h^\vee$”出自译本
      \texttt{witten.tex:1206--1208} 的 [译注]，不是原文。引用时写
      “据译本批注 \texttt{witten.tex:1206--1208} [译注]”，source 记
      \texttt{reconstruction}。
\end{enumerate}

\subsection{译本工程自述与实际的差异（三条）}

这三条不影响数学，但影响“能不能拿 README 当验收依据”，研讨班引用时须知：

\begin{enumerate}[leftmargin=1.8em,itemsep=2pt]
\item 交付说明称 \texttt{\% [Tutor Note: \dots]} 共 101 条；实际 \texttt{witten.tex}
      中 103 条、\texttt{校勘与审校批注汇总.md} 中 105 条。
\item 交付说明称“复核记录以 \texttt{\% ==== 复核记录 ====} 保留在 \texttt{witten.tex}”；
      实际 \texttt{witten.tex} 中 0 次，只在 \texttt{sections/*.tex}（14 次）与
      \texttt{build/tr7audit/witten\_fresh.tex}（12 次）中。
\item 交付说明称“凡不确定处批注‘存疑’”；实际 \texttt{witten.tex} 与校勘汇总里
      “存疑”均 0 次。\textbf{因此“未标存疑”不等于“已核对”}，坐标级核对仍以本附录
      与页图为准。
\end{enumerate}

\subsection{本次新发现（2026-11-13，S5/S6 证据官）}

\begin{enumerate}[leftmargin=1.8em,itemsep=2pt]
\item \texttt{sessions/s5/claims.jsonl} 的 12 条仍用旧 \texttt{status}
      （11 条 \texttt{accepted}、1 条 \texttt{provisional}）；CHARTER §5（行 146）只许
      \texttt{active/contested/refuted}。\texttt{check\_claims.py} 容忍旧值，故 0 错误
      掩盖了它。\texttt{s1--s4} 已归一，\texttt{s5} 漏改（详见
      \texttt{sessions/\_anchor\_report.md}）。
\item \texttt{c-s5-08} 的锚 “\texttt{[上野 p.~103 (5.29)(5.30)]}” 页码错：应为
      p.~111（见上第 10 条）。
\item \texttt{c-s2-10} 的锚 “\texttt{[src p.~385--388 §4.3]}”：§4.3 实际是
      p.~384--387（p.~388 已进入 §4.5），建议改为 p.~384--387。
\item \texttt{c-s6-07} 的锚 “\texttt{[原文 p.~351--399 全体]}” 无法核对，建议改成
      具体页或降级为 \texttt{open}。
\item \textbf{检查器盲区}：\texttt{check\_anchors.py} 只识别方括号
      “\texttt{[eq x.y]}”式锚点，s1 中裸露的 “eq 1.2/1.3/1.6/4.21/4.41” 与 s6 的
      “\texttt{[eq 3.5 p.373]}”未被检查；\texttt{check\_claims.py} 只核对以
      \texttt{[src}/\texttt{[原文} 开头的页锚，s1 中裸露的 “原文 p.$N$”（12 条）
      与全部教材页锚（河野/上野/Moore/李思）完全未核。因此“0 错误”只说明
      \textbf{已覆盖的那部分}一致。
\end{enumerate}

\subsection{引用格式}

正文引用本附录时写“见附录 \ref{app:errata}”；引用未核内容时写
\texttt{[待核]} 并注明待核对象（页码、版本或某教材某页），不要用“应该”“想必”。
